\section{Threat Analysis of DecMed}
\label{sec:threat-analysis}

\subsection{System Decomposition}

Before applying STRIDE, we decomposed DecMed into external dependencies,
entry points, assets, and trust levels. The decomposition identifies nine
external dependencies (IOTA network, Move smart contract package, PRE
Server, Gas Station, Redis, IPFS Cluster, Tauri client framework, BIP39
scheme, and PRE/AES-GCM cryptographic primitives) and 14 entry points
ranging from PRE Server REST endpoints to on-chain Move method invocations.

Key assets include: administrative and clinical EHR data, CapBAC capability
entries, BIP39 recovery phrases and derived private keys, PRE key material
(\texttt{kfrag}/\texttt{cfrag}), JWT access tokens, activation keys, and the
\texttt{GlobalAdminCap}/\texttt{ProxyCapEnum} capability objects---compromise
of the latter is equivalent to total privilege escalation.

Eight trust levels are defined, from unauthenticated actors (LK1) through
registered patients and healthcare personnel (LK2--LK5), Ministry of Health
administrators (LK6), up to privileged server-side components (LK7--LK8).
The resulting DFD captures all data flows among Patient Client, Hospital
Client, Ministry Client, PRE Server, Gas Station, IOTA smart contract, IPFS
Cluster, and Redis, with trust boundaries drawn at each component interface.
IOTA Network and Gas Station, as external trusted providers, are treated as
outside the threat model boundary; analysis is limited to their interactions
with internal components.

\subsection{STRIDE Identification}

Applying STRIDE per-interaction to the DFD yields \textbf{25 threats
(T1--T25)} distributed across all six categories, summarized in
Table~\ref{tab:stride-distribution}.

\begin{table}[htbp]
    \caption{STRIDE Threat Distribution in DecMed}
    \label{tab:stride-distribution}
    \begin{center}
        \begin{tabular}{l c p{4.8cm}}
            \toprule
            \textbf{Category} & \textbf{Count} & \textbf{Representative Threats} \\
            \midrule
            Spoofing           & 5 & Patient/personnel impersonation during authentication; Gas Station operator forgery \\
            Tampering          & 4 & In-transit event modification; clinical data tampering on IPFS; capability entry forgery \\
            Repudiation        & 5 & Denial of EHR access; denial of consent grant/revocation; denial of re-encryption \\
            Info.\ Disclosure  & 5 & Clinical data exposure in transit; JWT leakage; PRE key interception \\
            Denial of Service  & 3 & IPFS flooding; Gas Station exhaustion; ALS endpoint flooding \\
            Elev.\ of Privilege & 3 & JWT role escalation; \texttt{ProxyCapEnum} capture; activation key misuse \\
            \midrule
            \textbf{Total}     & \textbf{25} & \\
            \bottomrule
        \end{tabular}
    \end{center}
\end{table}

\subsection{Derivation of Audit Event Requirements}

For each threat, we determine the minimum evidence required to rebut a
repudiation claim or support an investigation. This yields 13 mandatory
audit event types (EV1--EV13), listed in Table~\ref{tab:event-catalog}.
Events are tied to the STRIDE categories they primarily address, ensuring
that coverage is risk-driven rather than ad hoc.

\begin{table}[htbp]
    \caption{Audit Event Catalog (EV1--EV13)}
    \label{tab:event-catalog}
    \begin{center}
        \begin{tabular}{c p{4cm} l}
            \toprule
            \textbf{ID} & \textbf{Event} & \textbf{Primary Category} \\
            \midrule
            EV1  & Patient authentication (success/fail)       & S \\
            EV2  & Healthcare personnel authentication         & S \\
            EV3  & Patient on-chain registration               & S, R \\
            EV4  & Access grant (consent creation)             & R \\
            EV5  & Access revocation                           & R \\
            EV6  & EHR read operation                         & R, ID \\
            EV7  & EHR create/update                          & T, R \\
            EV8  & Administrative metadata update              & T, R \\
            EV9  & Activation key issuance                    & EoP \\
            EV10 & Personnel registration via activation key  & EoP, S \\
            EV11 & PRE re-encryption execution                & R, ID \\
            EV12 & Policy violation / access denied            & EoP, S \\
            EV13 & ALS log rotation and archival               & T, DoS \\
            \bottomrule
            \multicolumn{3}{l}{\small S=Spoofing, T=Tampering, R=Repudiation,} \\
            \multicolumn{3}{l}{\small ID=Info.\ Disclosure, DoS=Denial of Service, EoP=Elev.\ of Privilege} \\
        \end{tabular}
    \end{center}
\end{table}

Each event carries a mandatory base schema (\texttt{record\_id},
\texttt{timestamp}, \texttt{source\_component}, \texttt{actor\_id},
\texttt{action}, \texttt{target\_object\_id}, \texttt{outcome},
\texttt{event\_signature}) plus event-specific evidence attributes. For
example, EV6 additionally records the IPFS CID of the retrieved clinical
data and the active capability entry reference, providing sufficient evidence
to attribute any specific read to an authenticated actor holding a valid
capability at that point in time.

We also identify five existing gaps in DecMed's current logging mechanisms
(M-TA-01 through M-TA-05), summarized in Table~\ref{tab:gaps}, which ALS
is designed to resolve.

\begin{table}[htbp]
    \caption{Identified Audit Logging Gaps in DecMed}
    \label{tab:gaps}
    \begin{center}
        \begin{tabular}{c p{5.5cm}}
            \toprule
            \textbf{ID} & \textbf{Gap Description} \\
            \midrule
            M-TA-01 & No centralized audit collector; logs scattered across components \\
            M-TA-02 & Narrow event coverage; only Move call executions recorded \\
            M-TA-03 & Audit records can be overwritten; non-append-only state \\
            M-TA-04 & No standardized audit record schema across components \\
            M-TA-05 & No retention strategy; evidence vulnerable to DLT data pruning \\
            \bottomrule
        \end{tabular}
    \end{center}
\end{table}